\pagenumbering{arabic}
\section{INTRODUCCIÓN}\label{cap:introduccion}

Al comienzo de cada año escolar, las familias colombianas reciben del colegio una lista de
útiles que deben conseguir completa y en pocos días. El gasto es alto y la presión de tiempo también: el 35\,\% de las familias destina entre \$300.000 y
\$500.000 por niño \citep{noauthor_resultados_2025} y el 70\,\% no planea la compra con anticipación
\citep{noauthor_resultados_2024}. Cuando la compra se retrasa o queda incompleta, quien
lo paga es el estudiante, que empieza el año sin los materiales que el colegio pidió para
sus clases.\par

Este proyecto estudió el problema de cerca en el
Colegio Agustiniano Ciudad Salitre. Una encuesta a sus acudientes, con 65 respuestas
recibidas y 60 analizables, mostró que el 65\,\% dedica cuatro horas o más a la compra, que el 68,3\,\% compra
en tres o más lugares y que la queja más frecuente es no encontrar el producto o que esté
agotado (43,3\,\%). Las entrevistas con directivos, docentes y acudientes confirmaron esa
fricción.

Las alternativas que existen hoy no resuelven ese recorrido completo: en las tiendas colombianas es el
acudiente quien debe traducir la lista del colegio en productos, y plataformas como TeacherLists, en Estados
Unidos, convierten la lista del salón en un carrito \citep{teacherlists_retail_2026}, pero, hasta donde se revisó (Capítulo~III), ninguna lee listas en papel o
en foto ni está disponible en Colombia.\par

Frente a este panorama, el proyecto construyó Tukit, un sistema de recomendación multiobjetivo con un
modelo B2B2C, en el que el colegio aporta la lista oficial y la familia hace la compra. El acudiente
elige el colegio y el grado de cada hijo, o sube su propia lista en imagen, PDF o texto libre, e indica
su presupuesto y sus preferencias. A cambio recibe hasta cuatro kits completos (Ahorrador, Premium, Fan y
A tu Medida) y puede cambiar cualquier producto antes de pagar (el pago es simulado en esta versión). Para
llegar a esos kits, el sistema recorre tres etapas.\par

\begin{itemize}
  \item \textbf{Percepción.} Un modelo multimodal de inteligencia artificial lee la lista y la convierte
        en renglones.
  \item \textbf{Normalización.} Cada renglón se clasifica (útil, texto escolar, licencia digital u otro),
        se separa la cantidad pedida del empaque en que se vende, y cada útil se busca en un catálogo de
        conceptos canónicos con una búsqueda híbrida: un canal léxico corrige la ortografía y exige que lo
        pedido esté nombrado, y un canal semántico reconoce sinónimos. Si el sistema no tiene certeza
        suficiente, el renglón pasa a revisión humana.
  \item \textbf{Optimización.} Un algoritmo calcula el frente de Pareto de los kits posibles, es decir,
        aquellos en los que no se puede bajar el costo sin ceder en la marca o en las preferencias. Con
        hasta 4.096 combinaciones las revisa todas; con más, usa NSGA-II \citep{deb_fast_2002}.
\end{itemize}

Esta estrategia toma de la generación aumentada por recuperación (RAG) \citep{lewis_retrieval-augmented_2020}
la indexación y la recuperación de productos, pero los kits los construye el optimizador, no un modelo de
lenguaje, y el equipo de Tukit los revisa antes de publicarlos (\emph{human-in-the-loop}). Todo corre sobre
cinco servicios de dominio con base propia y una puerta de enlace, es decir, seis unidades desplegables. El
diseño previó un séptimo servicio, la ingesta, que el sistema construido reemplazó por un catálogo curado.
El piloto se hizo en un solo colegio de Bogotá, y la encuesta describe a esa población sin generalizarse a
otros colegios; se eligió uno porque permitió llegar a familias de todos los niveles con la autorización del
rector y validar el flujo completo. Sumar otro colegio consiste en cargar y revisar su lista, sin cambiar el
sistema, pero la preferencia de sus familias habría que medirla de nuevo.\par

El trabajo siguió una metodología híbrida organizada en cinco fases, una por cada objetivo específico.
SCRUM \citep{schwaber_scrum_2020} organizó la construcción en ciclos de dos semanas; Design Thinking
\citep{brown_design_2008} guió la propuesta de valor; RUP \citep{kruchten_rup_2003} aportó la gestión de
riesgos y la trazabilidad; ATAM \citep{bass_software_2021} sirvió para elegir la arquitectura entre tres
alternativas; CRISP-DM \citep{chapman_crisp-dm_2000} fue el marco de los componentes que dependen de datos;
y los estándares IEEE 830 \citep{ieee_830_1998}, 1016 \citep{ieee_1016_2009} y 1058 \citep{ieee_1058_1998}
definieron la forma de los requerimientos (SRS), del diseño (SDD) y del plan del proyecto.\par

El documento sigue ese mismo orden, de entender el problema a construir la solución y luego a comprobarla.
El Capítulo~II presenta los objetivos y asigna cada fase (F1 a F5) a un objetivo (OE1 a OE5) y a una etapa
de CRISP-DM; el III reúne los conceptos y las soluciones existentes. El IV es la Fase~1 (OE1), con Design
Thinking y la comprensión del negocio: de las entrevistas y la encuesta salen los requerimientos según IEEE~830.
El V es la Fase~2 (OE2), donde ATAM elige la arquitectura y IEEE~1016 la documenta. El VI recoge las Fases~3
y~4 (OE3 y OE4) y recorre CRISP-DM, de la comprensión de los datos al despliegue, para la percepción, la
normalización y la optimización. El VII es la Fase~5 (OE5), con ISO/IEC~25010 y un protocolo registrado de
antemano para las pruebas y la validación con usuarios, y el VIII cierra con las conclusiones y el trabajo
futuro. Los anexos solo guardan material de reproducción y auditoría.\par
